\subsection{Upper bound for the cardinality of the fibres of an étale and separated map}

THEOREM 3. --- Let E and B be topological spaces and p: E → B an étale and separated map. Suppose that E is connected and that B is locally connected. Let W be a base of the topology of B. Then, for every point a of B, one has Card(E \(_{a}\) ) ≤ sup(Card(W), Card(N)).

Since the space B is locally connected, its topology has a base \(\mathcal{W}'\) consisting of connected open sets, for example the connected components of the open sets of \(\mathcal{W}\) (TG, I, p.~85, prop. 11). By virtue of Lemma 4 below, there exists a base \(\mathcal{V}\) of the topology of B composed of connected open sets and such that \(\mathrm{Card}(\mathcal{V}) \leqslant \mathrm{Card}(\mathcal{W})^2\) .

Lemma 4. --- Let B be a topological space and let W and W' be bases of the topology of B. There exists a subset V of W' which is a base of the topology of B and such that \(\text{Card}(\mathcal{V}) \leqslant \text{Card}(\mathcal{W})^{2}\) .

Let \(\mathcal{A} \subset \mathcal{W} \times \mathcal{W}\) be the set of pairs \((\mathrm{W}_1, \mathrm{W}_2)\) for which there exists \(\mathrm{W}' \in \mathcal{W}'\) such that \(\mathrm{W}_1 \subset \mathrm{W}' \subset \mathrm{W}_2\) ; let \(\varphi: \mathcal{A} \to \mathcal{W}'\) be a map such that one has \(\mathrm{W}_1 \subset \varphi(\mathrm{W}_1, \mathrm{W}_2) \subset \mathrm{W}_2\) for every pair \((\mathrm{W}_1, \mathrm{W}_2) \in \mathcal{A}\) , and let \(\mathcal{V} \subset \mathcal{W}'\) be the image of \(\varphi\) . We have

\[
\operatorname{Card} (\mathcal {V}) \leqslant \operatorname{Card} (\mathcal {A}) \leqslant \operatorname{Card} (\mathcal {W} \times \mathcal {W})
\]

(E, III, p.~25, prop. 3). Let us prove that V is a base of the topology of B. Let x be a point of B and U a neighbourhood of x. By hypothesis, x admits a neighbourhood \(W_{2} \in W\) contained in U, a neighbourhood \(W' \in W'\) contained in \(W_{2}\) and a neighbourhood \(W_{1} \in W\) contained in \(W'\) . One thus has \(W_{1} \subset W' \subset W_{2} \subset U\) . Then, \((W_{1}, W_{2}) \in \mathscr{A}\) and \(\varphi(W_{1}, W_{2})\) is a neighbourhood of x contained in U. According to prop. 3 of TG, I, p.~5, the set V is a base of the topology of B.

Let V be a base of the topology of B consisting of nonempty connected open sets and such that \(\text{Card}(\mathcal{V}) \leqslant \text{Card}(\mathcal{W})^{2}\) . Let U be the set of open sets U of E such that p induces a homeomorphism of U onto an open set V belonging to V. Every element of U is a nonempty connected open set. Since the map p is étale, according to definition 2 (I, p.~28) and prop. 3 of TG, I, p.~5, the set U is a base of the topology of E. Let us call a caterpillar every finite sequence \((\mathrm{U}_{0}, \ldots, \mathrm{U}_{n})\) of elements of U such that for \(1 \leqslant i \leqslant n\) , \(U_{i-1} \cap U_{i}\) is nonempty and \(p(\mathrm{U}_{i-1}) \cap p(\mathrm{U}_{i})\) is connected. Let S denote the set of finite sequences of elements of V and, for every caterpillar \(c = (\mathrm{U}_{0}, \ldots, \mathrm{U}_{n})\) , denote by \(p(c)\) the sequence \((p(\mathrm{U}_{0}), \ldots, p(\mathrm{U}_{n}))\) .

Lemma 5. --- a) Let \(c = (\mathrm{U}_0, \ldots, \mathrm{U}_n)\) and \(c' = (\mathrm{U}_0', \ldots, \mathrm{U}_m')\) be caterpillars such that \(\mathrm{U}_0 = \mathrm{U}_0'\) and \(p(c) = p(c')\) . Then, \(c = c'\) .

\begin{enumerate}
\def\labelenumi{\alph{enumi})}
\setcounter{enumi}{1}
\item
  Let U and U' be elements of U. There then exists a caterpillar \(c = (\mathrm{U}_{0}, \ldots, \mathrm{U}_{n})\) such that \(U_{0} = U\) and \(U_{n} = U'\) .
\item
  One necessarily has m = n. For every integer i such that \(0 \leqslant i \leqslant n\) , set \(V_{i} = p(U_{i})\) and denote by \(s_{i}\) and \(s_{i}^{\prime}\) the continuous sections of p over \(V_{i}\) with respective images \(U_{i}\) and \(U_{i}^{\prime}\) . To prove the equality \(c = c^{\prime}\) , we shall prove, by induction on i, the equality \(s_{i} = s_{i}^{\prime}\) for every integer i such that \(0 \leqslant i \leqslant n\) . By hypothesis, one has \(s_{0} = s_{0}^{\prime}\) . Let i be such that \(1 \leqslant i \leqslant n\) and \(s_{i-1} = s_{i-1}^{\prime}\) . Since \(U_{i-1} \cap U_{i}\) contains a point x, the continuous sections \(s_{i-1}\) and \(s_{i}\) coincide at \(p(x)\) , hence at every point of \(V_{i-1} \cap V_{i}\) , since it is connected (corollary 1 of I, p.~34). The same holds for \(s_{i-1}^{\prime}\) and \(s_{i}^{\prime}\) . For the same reason, \(s_{i}\) and \(s_{i}^{\prime}\) which coincide in \(V_{i-1} \cap V_{i}\) also coincide in \(V_{i}\) , whence the result.
\item
  If \(x\) and \(y\) are points of E, one says that a caterpillar \(c = (\mathrm{U}_0, \ldots, \mathrm{U}_n)\) connects \(x\) to \(y\) if \(x \in \mathrm{U}_0\) and \(y \in \mathrm{U}_n\). In the set E, let R be the relation “there exists a caterpillar that connects \(x\) to \(y\) ”. The relation R is reflexive since \(\mathcal{U}\) is a covering of E. It is evidently symmetric. Let us prove that it is transitive: let \(x, y, z\) be three points of E, \((\mathrm{U}_0, \ldots, \mathrm{U}_n)\) and \((\mathrm{U}_0', \ldots, \mathrm{U}_m')\) are caterpillars connecting \(x\) to \(y\) and \(y\) to \(z\) respectively. Let \(\mathrm{U} \in \mathcal{U}\) be a neighborhood of \(y\) contained in \(\mathrm{U}_n \cap \mathrm{U}_0'\) ; we have \(p(\mathrm{U}_n) \cap p(\mathrm{U}) = p(\mathrm{U}_0') \cap p(\mathrm{U}) = p(\mathrm{U})\) , and, since U is connected, the sequence \((\mathrm{U}_0, \ldots, \mathrm{U}_n, \mathrm{U}, \mathrm{U}_0', \ldots, \mathrm{U}_m')\) is a caterpillar that connects \(x\) to \(z\) . The equivalence classes under R are open, hence also closed. Since E is connected, it follows that any two points of E are always connected by a caterpillar.
\end{enumerate}

Let U and U′ be elements of U, let x be a point of U, and let \(x'\) be a point of U'. There exists a caterpillar \((\mathrm{U}_{2},\ldots,\mathrm{U}_{n-2})\) connecting x to \(x'\) . Let \(U_{1}\) and \(U_{n-1}\) be open subsets of U such that \(x \in U_{1}\) , \(x' \in U_{n-1}\) , \(U_{1} \subset U \cap U_{2}\) , \(U_{n-1} \subset U' \cap U_{n-2}\) . Set \(U_{0} = U\) , \(U_{n} = U'\) . Then the sequence \((\mathrm{U}_{0},\ldots,\mathrm{U}_{n})\) is a caterpillar.

Let us now prove the theorem. Let U be an element of \(\mathcal{U}\) and let C be the set of caterpillars \((\mathrm{U}_0,\dots ,\mathrm{U}_n)\) such that \(\mathrm{U}_0 = \mathrm{U}\). The map from C to S which assigns to a caterpillar \(c = (\mathrm{U}_0,\dots ,\mathrm{U}_n)\) of C the sequence \(p(c) = (p(\mathrm{U}_{0}), \ldots, p(\mathrm{U}_{n}))\) is injective (Lemma 5, (a)), hence (1) \(\text{Card}(\mathrm{C}) \leqslant \text{Card}(\mathrm{S}).\)

The map from C to \(\mathcal{U}\) which, to \(c = (\mathrm{U}_0,\dots ,\mathrm{U}_n)\in \mathrm{C}\) , associates the open set \(\mathrm{U}_n\) is surjective by the second part of Lemma 5, hence

\[
\operatorname{Card} (\mathcal {U}) \leqslant \operatorname{Card} (\mathrm{C}).\tag{2}
\]

For every point \(x\) of E, choose an open set \(\mathrm{U}_x\) of \(\mathcal{U}\) such that \(x \in \mathrm{U}_x\) . If \(x\) and \(y\) are distinct points of the same fiber \(\mathrm{E}_a\) of \(p\) , we have \(\mathrm{U}_x \neq \mathrm{U}_y\) since \(p \mid \mathrm{U}_x\) is injective. We therefore have, for \(a \in \mathrm{B}\) ,

\[
\operatorname{Card} \left(\mathrm{E} _ {a}\right) \leqslant \operatorname{Card} (\mathcal {U}).\tag{3}
\]

Finally, if V is a finite set, the set S of finite sequences of elements of V is countable (E, III, p.~49, prop. 1), hence

\[
\operatorname{Card} (\mathrm{S}) \leqslant \operatorname{Card} (\mathbf {N}).\tag{4}
\]

If \(\mathcal{V}\) is infinite, we have

\[
\operatorname{Card} (\mathrm{S}) = \operatorname{Card} (\mathcal {V}) \leqslant \operatorname{Card} (\mathcal {W}) ^ {2} = \operatorname{Card} (\mathcal {W})\tag{5}
\]

by the corollary of E, III, p.~50 and Corollary 1 of E, III, p.~49. The theorem follows from relations (1) to (5) above.

Remark. --- By the preceding, if the topology of B has a countable base, then so does that of E, and the fibers of E are countable (cf. TG, I, p.~88, Poincaré-Volterra theorem).

